\section{Stage 1 methodology}\label{sec:stage1}

Stage~1 consists of three steps: setting up a verifiable environment
(\cref{sec:env}), selecting the video (\cref{sec:selection}),
and preprocessing it into frozen, hashed artifacts
(\cref{sec:preprocessing}). All steps are implemented in the \memvid{} Python
package and exposed as \cmd{memvid <command>}; the commands are named in
the text. Parameter values quoted below are those of
\cmd{configs/default.yaml} and are listed in \cref{tab:stage1-params}, which is
generated from the configuration rather than copied.

\subsection{Step 1: environment and reproducibility}\label{sec:env}

\paragraph{Software.}
The package targets Python $\ge 3.10$ and uses a \cmd{src} layout with
dependencies declared in \cmd{pyproject.toml}: a small core (pydantic,
PyYAML, NumPy, Pillow, tqdm) and optional groups for speech recognition
(faster-whisper), the models (PyTorch, transformers), sentence embeddings,
figures, the demo and development tools. ffmpeg and ffprobe are system
binaries.

\paragraph{Environment check.}
\cmd{memvid check-env} records the Python and OS versions, the ffmpeg version,
the version of every relevant package, and the GPU, and fails if a hard
requirement is missing. The report is stored as JSON and exported to the
environment table below.

\paragraph{Reproducibility measures.}
\begin{enumerate}[nosep]
  \item \emph{Typed configuration.} Every run is defined by one YAML file
        validated against a schema that rejects unknown keys, so a misspelled
        parameter fails instead of silently taking its default.
  \item \emph{Determinism and freezing.} A single seed controls every random
        choice of the package (the dev/test split and the audit samples).
        Speech recognition is not fully deterministic under its default
        temperature fallback, which re-decodes difficult windows by sampling
        (configurable; disabling it gives deterministic decoding at some cost
        in quality). Identity of the inputs to later stages is therefore
        guaranteed by freezing the transcript, not by re-running ASR.
  \item \emph{Provenance manifest.} For every derived artifact the manifest
        stores its SHA-256, the producing step, a hash of the configuration
        sections that determined it, and the package version
        (\cref{sec:provenance}).
  \item \emph{Pinned versions.} The exact environment is recorded with
        \cmd{pip freeze}, and the commit used for reported results is tagged.
\end{enumerate}

\generated{environment.tex}{environment table, produced by
\cmd{memvid check-env} followed by \cmd{memvid export-tex}.}

\subsection{Step 2: video selection}\label{sec:selection}

The choice of video determines whether the research questions can be answered
at all: a video without recurring entities cannot test entity resolution, and
one without state changes cannot test updating. The video should therefore
last at least 20 minutes, as the brief requires (longer videos make evidence
age, and thus RQ2, observable over a wider range), have an audio track with
speech, contain objects and people that reappear in distant parts of the
video, and show objects whose state changes (painted, filled, assembled).

The video is placed in \cmd{data/\allowbreak raw/\allowbreak } and probed
(\cmd{memvid probe}): duration, resolution, frame rate, codecs, presence of an
audio stream, and the SHA-256 of the file, which identifies the exact source
for all later results. A video shorter than 20 minutes is reported with a
warning. The questions (Stage~2) are written before any memory exists, so they
cannot be tailored to what a particular organisation happens to capture.

\subsection{Step 3: preprocessing}\label{sec:preprocessing}

\subsubsection{Audio extraction}
The first audio stream is decoded to 16-bit PCM, mono, at 16~kHz, the native
input format of Whisper (\cmd{memvid extract-audio}). The duration of the
extracted audio is compared with the video duration; a difference above
0.5~s aborts the step, because every transcript timestamp is expressed in
audio time and would otherwise be misaligned with the frames.

\subsubsection{Speech recognition}\label{sec:asr}
Transcription uses faster-whisper~\citep{fasterwhisper} with the
\texttt{large-v3} model~\citep{radford2023whisper} and beam size~5
(\cmd{memvid transcribe}). Two settings deviate from library defaults, both
motivated by known failure modes on long recordings:
\begin{itemize}[nosep]
  \item \emph{Voice-activity filtering} (\texttt{vad\_filter=true}): non-speech
        regions are removed before decoding, which suppresses fluent but
        hallucinated text on silence and background
        sounds~\citep{koenecke2024careless}.
  \item \emph{No conditioning on previous text}
        (\texttt{condition\_on\_previous\_text=false}): with conditioning, one
        decoding error is fed back as the prompt for the next window and can
        start a repetition loop that persists for minutes.
\end{itemize}
The output is a list of segments $s_j=[a_j,b_j]$ with text, average
log-probability and no-speech probability. Segments whose no-speech
probability exceeds 0.6 are counted as suspicious and reported, not removed:
removal would be an unverified editorial decision on the evidence. Segment
timestamps are not checked by hand: transcripts enter the memory at chunk
granularity ($L=20$~s), where a boundary error of about a second changes at
most which of two adjacent chunks a segment is attached to.

\subsubsection{Chunking}\label{sec:chunking}
With chunk length $L$ and $n_0=\lfloor T/L\rfloor$, the base chunks are
$[iL,(i+1)L)$ for $i=0,\dots,n_0-1$. The remainder $r=T-n_0L$ forms an
additional chunk $[n_0L,T)$ if $r\ge\alpha L$, and is otherwise merged into
the last chunk. Hence every chunk has length in $[\alpha L,(1+\alpha)L)$
unless $T<L$. Defaults: $L=20$~s, $\alpha=0.5$ (\cmd{memvid chunk}).

The choice of $L$ is a trade-off. Smaller chunks give finer provenance, since
$L$ is the temporal resolution of every grounding result, but increase the
number of perception calls ($n\approx 180$ for one hour at $L=20$~s) and
fragment events that last longer than a chunk across several observations,
which the memory must then reassemble. A sensitivity check with
$L\in\{10,30\}$~s is planned as an ablation once perception exists.

\paragraph{Transcript assignment.}
Segment $s_j$ is attached to chunk $c_i$ iff
\begin{equation}
  \operatorname{ov}(s_j,c_i) = \min(b_j,v_i)-\max(a_j,u_i) \;\ge\;
  \min\bigl(\delta,\; b_j-a_j\bigr),
  \qquad \delta=0.25\text{ s}.
\end{equation}
The $\min$ with the segment length ensures that very short segments are
attached to the chunk containing them. A segment that straddles a boundary by
more than $\delta$ on both sides is attached to both chunks. This duplication
is deliberate: speech near a boundary is evidence for both chunks, and
assigning it to only one would make the observation of the other chunk
incomplete. A segment overlapping a chunk by less than $\delta$ is a
timestamp-noise artefact and is not attached to it.

\subsubsection{Frame sampling}
Frame $k$ of chunk $[u,v)$ is requested at the midpoint of the $k$-th of $N$
equal sub-intervals,
\begin{equation}
  t_k = u + \bigl(k+\tfrac12\bigr)\frac{v-u}{N}, \qquad k=0,\dots,N-1,
\end{equation}
(\cmd{memvid sample-frames}). Unlike endpoint sampling, midpoints never fall on
a chunk boundary, so adjacent chunks never share a frame, and the spacing
$(v-u)/N$ stays uniform across boundaries. With $N=8$ and $L=20$~s the
effective rate is $0.4$ frames/s. This rate is a real limitation and is stated
as such: a purely visual event shorter than $L/N=2.5$~s can fall between two
samples and be invisible to perception. It bounds perception recall and must
be kept in mind when interpreting perception misses
(\cref{sec:attribution}).

\paragraph{Decoding and timestamp accuracy.}
Frames are extracted by calling ffmpeg once per chunk, seeking on the input
(frame-accurate when decoding) and applying a \texttt{select} filter that keeps,
for each target $t_k$, the \emph{first frame whose presentation time is at or
after} $t_k$. The \texttt{showinfo} filter reports the true presentation time
$t_k^{\mathrm{act}}$ of each kept frame, which is stored alongside the request.
The selection rule bounds the timing error to $0\le t_k^{\mathrm{act}}-t_k<1/f$
for frame rate $f$. Frames are resized so that the longer side is at most
768~px (the visual-token cost of the VLM grows with resolution) and stored as
JPEG at quality 90.

\paragraph{Verification of the frame sampler.}
The sampler was verified on a synthetic test pattern (ffmpeg
\texttt{testsrc2}, 25~fps), in which every frame differs from every other.
Each extracted frame was compared with all frames obtained by an independent
sequential decoder (OpenCV): in every case the extracted image was
pixel-identical to the frame at the reported $t_k^{\mathrm{act}}$, and the
timing error lay in $[0,1/f)$, including targets within one frame of the end
of the stream.

\subsubsection{Provenance and freezing}\label{sec:provenance}
All artifacts of a video are written atomically (temporary file and rename)
into \cmd{data/\allowbreak derived/\allowbreak<video\_id>/} and registered in a manifest with their
hashes (\cref{app:formats}). Frames are registered as one tree hash over the
sorted pairs of relative path and file hash. After preprocessing,
\cmd{memvid freeze} verifies all hashes and marks the derived data as frozen.
From then on, any step that would overwrite a frozen artifact refuses to run
unless explicitly forced, and \cmd{memvid verify} detects any modification.
Every later stage records the manifest hash it consumed, so a result can
always be traced to the exact inputs that produced it.

\generated{stage1_params.tex}{parameter table, produced by
\cmd{memvid export-tex}.}

\subsection{Software verification}
During development the package was checked by an automated test suite (75
tests, all passing): timestamp parsing, configuration validation, chunk
construction and its edge cases, the transcript assignment rule, frame timing
and end-of-stream handling, the memory invariants (budget never exceeded,
provenance preserved, log replay reproduces the memory), scoring, and the
complete pipeline on a synthetic video. The submitted repository keeps the
smoke test instead (README), which runs every stage end to end without a GPU.
